\documentclass[12pt]{article}
\usepackage[utf8]{inputenc}
\usepackage[T1]{fontenc}
\usepackage{hyperref}
\title{Pumi, Southern Language Resources for K-12 Curriculum Development}
\author{Compiled by Claude Sonnet 5.5 (claude-sonnet-5-5) for Andr\'as Kornai}
\date{}
\begin{document}
\maketitle

\section{Introduction}
Southern Pumi (ISO 639-3 pmj; Glottocode sout2729) is a Sino-Tibetan language of Yunnan, China. Omniglot and the Endangered Languages Catalogue give about 22,000 speakers in Lanping, Ludian and Ninglang counties. The Catalogue lists it as vulnerable. The searches found no school materials. Descriptions of Pumi (Prinmi) in the literature often cover Central or Western dialects, and the sources disagree on how the dialects line up with the two ISO codes. This is flagged below.

\section{Primary Educational Resources}
No primer, textbook or school programme for Southern Pumi was found. A pinyin-based Roman spelling exists but is not widely used, and some teaching in Ninglang has used the Tibetan script (Wikipedia, Omniglot).

\section{Linguistic and Reference Materials}
\begin{itemize}
\item Endangered Languages Catalogue, Southern Pumi bibliography (lists Ding's \emph{A Grammar of Prinmi}, Brill 2014, which is based on the Central dialect): \url{https://endangeredlanguages.com/lang/5608/bibliography}
\item Omniglot page on Pumi: \url{https://omniglot.com/writing/pumi.htm}
\item Wikipedia article on the Pumi language: \url{https://en.wikipedia.org/wiki/Pumi_language}
\item HKU Scholars Hub paper on Ding's phonology-based dialect division of Prinmi: \url{https://hub.hku.hk/handle/10722/188206}
\end{itemize}
No dictionary specific to Southern Pumi was found.

\section{Religious and Community Resources}
No religious or community language resources for Southern Pumi were found.

\section{Digital and Audio Resources}
No audio, video, corpus or software was found.

\section{Curriculum Development Context}
Young people are thought to be using the language as a first language in decreasing numbers, so there is room for early intervention. Everything a curriculum needs (orthography decision, recordings, texts) would have to be created. The Prinmi grammar literature is the main starting point, with care taken over which dialect it describes.
\end{document}